\documentclass[12pt]{article}
\usepackage[utf8]{inputenc}
\usepackage[T1]{fontenc}
\usepackage[margin=1in]{geometry}
\usepackage{setspace}
\usepackage{xcolor}
\usepackage[colorlinks=true, urlcolor=blue!60!black, linkcolor=blue!60!black]{hyperref}
\setlength{\parindent}{0em}
\setlength{\parskip}{0.45em}
\renewcommand{\baselinestretch}{1.0}
\begin{document}
\begin{center}
{\Large \textit{Geduldig v. Aiello}, 417 U.S. 484 (1974)\par}
\vspace{0.3em}
{\normalsize Case Brief --- One-Page Edition\par}
\end{center}

\textbf{Memory Jogger.} Excluding normal pregnancy from state disability insurance is not sex discrimination: it divides ``pregnant women and nonpregnant persons'' (\href{https://supreme.justia.com/cases/federal/us/417/484/}{Justia}).

\textbf{Facts.} California's self-supporting, employee-funded disability insurance program excluded disabilities from pregnancy. After the State paid abnormal-pregnancy claims, only Jacqueline Jaramillo's normal-pregnancy claim remained.

\textbf{Procedural History.} A three-judge district court held the exclusion irrational; direct appeal.

\textbf{Votes.} Reversed 6--3: Stewart for the Court; Brennan, with Douglas and Marshall, dissented.

\textbf{Holding.} A pregnancy classification is not a sex classification absent proof of pretext, so rational basis review applies and the exclusion stands.

\textbf{Reasoning.} Covering pregnancy would raise costs by about a third and end the fund's self-support; a State may address a problem one step at a time. The program insured the same risks for men and women. Brennan answered that the fund covered male-only conditions such as prostatectomies, creating a double standard.

\textbf{Examples.} \emph{Same-Side:} \textit{Skrmetti}'s gender-dysphoria treatment ban. \emph{Opposite-Side:} refusing to hire previously pregnant women. \emph{Fence-Sitter:} a contraceptive exclusion alongside erectile-dysfunction coverage.

\textbf{Critique.} Pregnancy is a perfect proxy for sex, and Congress rejected the Title VII sequel in the Pregnancy Discrimination Act.

\textbf{Key Quotation.} ``Absent a showing that distinctions involving pregnancy are mere pretexts designed to effect an invidious discrimination against the members of one sex or the other, lawmakers are constitutionally free to include or exclude pregnancy from the coverage of legislation such as this on any reasonable basis, just as with respect to any other physical condition.'' --- Justice Stewart, 417 U.S. at 496--497 n. 20.

\end{document}
